% Physical pages: 441--443 (printed pages 418--420).
% Figures: none.
% Uncertainties: none.

\subsection[Varieties of Statistics]{\texorpdfstring{Varieties of Statistics\protect\footnotemark}
 {Varieties of Statistics}}
\footnotetext{This section lies somewhat out of the book's main line of
development, and may be omitted in a first reading.}

We can now take up a question raised in Chapter 4: what are the possibilities
for the change of state-vectors when we interchange identical particles?

For this purpose, we will consider the preparation of the initial or final
states in a scattering process. Suppose that a set of indistinguishable
particles in either of these states is brought to a particular configuration
with momenta $\mathbf p_1,\mathbf p_2$, etc. from a standard configuration with
momenta $\mathbf k_1,\mathbf k_2$, etc., by some sort of slowly varying external
fields, keeping the particles far enough apart in the process to justify the
use of non-relativistic quantum mechanics. (Spin indices are not shown
explicitly here; they should be understood to accompany momentum labels.) To
calculate the amplitude for this process we can use the path-integral
method,\footnote{Here I am following the discussion of Laidlaw and C.
DeWitt~\hyperref[ref:9.11]{[11]}, except that they apply the path-integral method to the whole
scattering process, rather than just the preparation of initial or final
states. In a relativistic theory the possibility of particle creation and
annihilation makes it necessary to apply the path-integral method to fields
rather than particle orbits. For us, this is not a problem, because we limit
such calculations to sufficiently early or late times, when the particles
participating in a scattering process are all far apart.} taking the $q$s and
$p$s of Section 9.1 as particle positions and momenta, rather than fields and
their canonical conjugates. These always satisfy canonical commutation rather
than anticommutation relations, whether or not the particles are bosons or
fermions or something else, so at this point we are not committing ourselves
to any particular statistics. The path-integral formula \eqref{eq:9.1.34} gives an
amplitude
$\langle\mathbf p_1,\mathbf p_2,\ldots\mid
\mathbf k_1,\mathbf k_2,\ldots\rangle_D$ as an integral over paths in which one
particle is brought continuously from momentum $\mathbf k_1$ to momentum
$\mathbf p_1$, another identical particle is brought continuously from
momentum $\mathbf k_2$ to momentum $\mathbf p_2$, and so on. The subscript `D'
indicates that this is the amplitude we would calculate for distinguishable
particles. In particular, this amplitude is symmetric under permutations of
the $\mathbf p$s \emph{and} simultaneous permutations of the $\mathbf k$s, but
has no particular symmetry under separate permutations of the $\mathbf p$s or
$\mathbf k$s. But if the particles are really indistinguishable, then there are
other paths that are topologically distinct, but that yield the same final
configuration. For space dimensionality $d\geq3$, the only such
paths\footnote{This is expressed formally in the statement that the first
homotopy group of configuration space in $d\geq3$ is the permutation
group~\hyperref[ref:9.12]{[12]}. By `configuration space' for $N$ distinguishable particles is
meant the space of $N$ $d$-vectors, excluding $d$-vectors that coincide with
(or are within an arbitrary limiting distance of) each other, and identifying
configurations that differ only by a permutation of the vectors.} are those
that take $\mathbf k_1,\mathbf k_2,\ldots$ into some non-trivial permutation
$\mathcal P$ of $\mathbf p_1,\mathbf p_2,\ldots$. Hence the true amplitude
should be written
\begin{equation}
\langle\mathbf p_1,\mathbf p_2,\ldots\mid
 \mathbf k_1,\mathbf k_2,\ldots\rangle
=\sum_{\mathcal P}C_{\mathcal P}
 \langle\mathbf p_{\mathcal P1},\mathbf p_{\mathcal P2},\ldots\mid
 \mathbf k_1,\mathbf k_2,\ldots\rangle_D,
\tag{9.7.1}\label{eq:9.7.1}
\end{equation}
the sum running over all $N!$ permutations of the $N$ indistinguishable
particles in the state, and $C_{\mathcal P}$ a set of complex constants. These
amplitudes must satisfy a composition rule appropriate for indistinguishable
particles:
\begin{align}
&\langle\mathbf p_1,\mathbf p_2,\ldots\mid
 \mathbf k_1,\mathbf k_2,\ldots\rangle
\notag\\
&\quad=\frac{1}{N!}\int d^3q_1\,d^3q_2\cdots
 \langle\mathbf p_1,\mathbf p_2,\ldots\mid
 \mathbf q_1,\mathbf q_2,\ldots\rangle
 \langle\mathbf q_1,\mathbf q_2,\ldots\mid
 \mathbf k_1,\mathbf k_2,\ldots\rangle.
\tag{9.7.2}\label{eq:9.7.2}
\end{align}
Using Eq. \eqref{eq:9.7.1}, this is the requirement
\begin{align*}
&\sum_{\mathcal P}C_{\mathcal P}
 \langle\mathbf p_{\mathcal P1},\mathbf p_{\mathcal P2},\ldots\mid
 \mathbf k_1,\mathbf k_2,\ldots\rangle_D
\\
&\quad=\frac{1}{N!}\sum_{\mathcal P',\mathcal P''}
 C_{\mathcal P'}C_{\mathcal P''}\int d^3q_1\,d^3q_2\cdots
 \langle\mathbf p_{\mathcal P'1},\mathbf p_{\mathcal P'2},\ldots\mid
 \mathbf q_1,\mathbf q_2,\ldots\rangle_D
\\
&\qquad\cdot
 \langle\mathbf q_{\mathcal P''1},\mathbf q_{\mathcal P''2},\ldots\mid
 \mathbf k_1,\mathbf k_2,\ldots\rangle_D.
\end{align*}
Applying a permutation $\mathcal P''$ to both the initial and final states in
the first amplitude on the right, this is
\begin{align*}
&\sum_{\mathcal P}C_{\mathcal P}
 \langle\mathbf p_{\mathcal P1},\mathbf p_{\mathcal P2},\ldots\mid
 \mathbf k_1,\mathbf k_2,\ldots\rangle_D
\\
&\quad=\frac{1}{N!}\sum_{\mathcal P',\mathcal P''}
 C_{\mathcal P'}C_{\mathcal P''}\int d^3q_1\,d^3q_2\cdots
 \langle\mathbf p_{\mathcal P''\mathcal P'1},
 \mathbf p_{\mathcal P''\mathcal P'2},\ldots\mid
 \mathbf q_{\mathcal P''1},\mathbf q_{\mathcal P''2},\ldots\rangle_D
\\
&\qquad\cdot
 \langle\mathbf q_{\mathcal P''1},\mathbf q_{\mathcal P''2},\ldots\mid
 \mathbf k_1,\mathbf k_2,\ldots\rangle_D.
\end{align*}
But the amplitudes
$\langle\mathbf p_1,\mathbf p_2,\ldots\mid
\mathbf k_1,\mathbf k_2,\ldots\rangle_D$ satisfy the composition rule for
distinguishable particles
\begin{align}
&\langle\mathbf p_1,\mathbf p_2,\ldots\mid
 \mathbf k_1,\mathbf k_2,\ldots\rangle_D
\notag\\
&\quad=\int d^3q_1\,d^3q_2\cdots
 \langle\mathbf p_1,\mathbf p_2,\ldots\mid
 \mathbf q_1,\mathbf q_2,\ldots\rangle_D
 \langle\mathbf q_1,\mathbf q_2,\ldots\mid
 \mathbf k_1,\mathbf k_2,\ldots\rangle_D,
\tag{9.7.3}\label{eq:9.7.3}
\end{align}
so the composition rule for the physical amplitudes may be written
\begin{align*}
&\sum_{\mathcal P}C_{\mathcal P}
 \langle\mathbf p_{\mathcal P1},\mathbf p_{\mathcal P2},\ldots\mid
 \mathbf k_1,\mathbf k_2,\ldots\rangle_D
\\
&\quad=\frac{1}{N!}\sum_{\mathcal P',\mathcal P''}
 C_{\mathcal P'}C_{\mathcal P''}
 \langle\mathbf p_{\mathcal P''\mathcal P'1},
 \mathbf p_{\mathcal P''\mathcal P'2},\ldots\mid
 \mathbf k_1,\mathbf k_2,\ldots\rangle_D,
\end{align*}
which will be satisfied if and only if
\begin{equation}
C_{\mathcal P'\mathcal P''}=C_{\mathcal P'}C_{\mathcal P''}.
\tag{9.7.4}\label{eq:9.7.4}
\end{equation}
That is, the coefficients $C_{\mathcal P}$ must furnish a one-dimensional
representation of the permutation group. But the permutation group has only
two such representations: one is the identity, with $C_{\mathcal P}=+1$ for
all permutations, and the other is the alternating representation, with
$C_{\mathcal P}=+1$ or $C_{\mathcal P}=-1$ according to whether $\mathcal P$ is
an even or odd permutation. These two possibilities correspond to Bose or
Fermi statistics, respectively.\footnote{There has been much discussion in the
literature of possibilities other than Bose or Fermi statistics, often under
the label \emph{parastatistics}. It has been shown~\hyperref[ref:9.13]{[13]} that parastatistics
theories in $d\geq3$ space dimensions are equivalent to theories in which all
particles are ordinary fermions or bosons, but carrying an extra quantum
number, so that wave functions could have unusual properties under
permutations of momenta and spins.}

The nice feature of this argument is that it makes it clear why the case of two
space dimensions is an exception. In this case there is a much richer variety
of topologically distinct paths.\footnote{This is expressed in the statement
that the first homotopy group of configuration space in two space dimensions
is not the permutation group, but a larger group known as the braid
group~\hyperref[ref:9.14]{[14]}.} For instance, a path in which one particle circles another a
definite number of times cannot be deformed into a path where it does not. In
consequence, in two space dimensions it is possible to have
\emph{anyons}~\hyperref[ref:9.15]{[15]}, particles with more general permutation properties than
just Fermi or Bose statistics~\hyperref[ref:9.8a]{[8a]}.
